% ============================================================
% §2  The Phi field — sources, superposition, interference
% ============================================================
\section{The space-generating field}\label{sec:pole}

\subsection{Definition of the field $\Phi$}

\begin{definition}[Spatial field $\Phi$]\label{def:Phi}
$\Phi$ is a real scalar field whose value at a given point measures
the \textbf{density of generated space} at that location.  It is not
a field on a ready-made manifold --- it is the field that
\emph{constitutes} the manifold:
\begin{itemize}[nosep]
  \item $\Phi = 0$: no space (nothingness $\Nzero$),
  \item $\Phi > 0$: space exists, its ``density'' equals~$\Phi$,
  \item $\Phi \to \infty$: extremely dense space
        (interior of massive objects).
\end{itemize}
\end{definition}

$\Phi$ is not an \textbf{additional field} in space.
It is the field \textbf{from which} space is composed.
Metric, distance, causality --- these are properties of the
configuration of~$\Phi$, not separate entities.

\begin{remark}[Physical interpretation]
One may think of $\Phi$ as a ``density of aether'' --- but not of a
material medium; rather, of the very \emph{capacity for spatial
relations to exist}.  Where $\Phi = 0$ there is nothing --- not even
emptiness.  Where $\Phi > 0$, objects can coexist, interact, and move
relative to one another.
\end{remark}

% ----------------------------------------------------------
\subsection{Matter as a source of $\Phi$}\label{ssec:zrodla}

Every fundamental particle generates a contribution to the
field~$\Phi$.  This is not an ``external'' coupling (matter in space
coupled to a field); it is a \textbf{constitutive act}: a particle
\emph{is} a local excitation that simultaneously produces the spatial
field around itself.

\begin{axiom}[Matter generates space]\label{ax:zrodlo}
Each source of mass-energy $\rho_a$ produces a contribution to~$\Phi$:
\begin{equation}\label{eq:Phi-source}
    \mathcal{D}[\Phi_a] = q_a \cdot \rho_a,
\end{equation}
where $q_a > 0$ is the \textbf{spatial generation constant} of
particle type~$a$ (analogous to charge, but for space), and
$\mathcal{D}$ is the operator governing the dynamics of~$\Phi$.
\end{axiom}

\begin{remark}[Generation constant $q_a$]
The constant $q_a$ specifies how strongly a given particle type
generates space.  It may depend on mass, spin, or other particle
properties.  It is a new physical parameter, intrinsic to TGP and
absent from standard physics.  In the homogeneous limit ($q_a = q$
for all types) the theory reduces to a one-parameter form.
\end{remark}

\begin{remark}[Nature of the operator $\mathcal{D}$]\label{rem:D-natura}
The operator $\mathcal{D}$ in Eq.~\eqref{eq:Phi-source} is not a
standard Laplacian on a pre-existing manifold --- that would violate
Axiom~\ref{ax:przestrzen}.  It must be \textbf{internal to $\Phi$}:
defined by the very configuration of the field that constitutes
space.  Its explicit form is derived in Section~\ref{sec:formalizm}
(Definition~\ref{def:D}).

For reference, the linear part of the operator reads
\begin{equation}\label{eq:D-preview}
  \mathcal{D}[\Phi]\big|_{\mathrm{lin}}
  = \frac{1}{\PhiZero}\,\nabla^2\Phi,
\end{equation}
while the full field equation includes the nonlinear
self-interference term $\mathcal{N}[\Phi]$
(Definition~\ref{def:N}, Section~\ref{sec:formalizm}):
\begin{equation}\label{eq:N-preview}
  \mathcal{N}[\Phi]
  = \frac{\alpha}{\PhiZero}\,\frac{(\nabla\Phi)^2}{\Phi}
  + \beta\,\frac{\Phi^2}{\PhiZero^2}
  - \gamma\,\frac{\Phi^3}{\PhiZero^3},
\end{equation}
where $\alpha = 2$ (from variation of the action),
$\beta, \gamma > 0$, $[\beta] = [\gamma] = \mathrm{L}^{-2}$.
\end{remark}

% ----------------------------------------------------------
\subsection{Single source: spherical generated space}%
\label{ssec:pojedyncze}

For an isolated, static point particle of mass~$M$ located
at~$\mathbf{x}_0$, the generated field has spherical symmetry:
\begin{equation}\label{eq:Phi-single}
    \Phi_{\mathrm{single}}(r)
    = \frac{q\,M}{4\pi}\,\phi(r),
    \qquad
    r = |\mathbf{x} - \mathbf{x}_0|,
\end{equation}
where $\phi(r)$ is the \textbf{generation profile} --- a function of
distance from the source.

The profile $\phi(r)$ is \textbf{not} assumed a priori to be
$e^{-mr}/r$ (Yukawa) or $1/r$ (Coulomb).  Its form follows from the
full, nonlinear dynamics of~$\Phi$ and must be \textbf{derived} from
the equations of the theory.  At the present stage one may only
assume general properties:
\begin{enumerate}[nosep]
  \item $\phi(r) > 0$ everywhere (space is generated, never
        ``negative''),
  \item $\phi(r) \to \infty$ as $r \to 0$ (spatial density diverges
        at the source centre),
  \item $\phi(r) \to 0$ as $r \to \infty$ (space vanishes far from
        the source),
  \item $\phi(r)$ is non-monotonic --- it possesses an inflection
        point or local extremum, giving rise to three regimes
        (Section~\ref{sec:rezimy}).
\end{enumerate}

% ----------------------------------------------------------
\subsection{Multiple sources: superposition and interference}%
\label{ssec:interferencja}

In the presence of many particles the total field~$\Phi$ is not a
simple sum of individual contributions.  It is the result of
\textbf{nonlinear interference}.

\subsubsection{Linear approximation (weak field)}

In the weak-field limit ($\Phi \ll \Phi_{\mathrm{char}}$, where
$\Phi_{\mathrm{char}}$ is the nonlinearity scale) the contributions
superpose linearly:
\begin{equation}\label{eq:Phi-linear}
    \Phi(\mathbf{x})
    \simeq \sum_{i} \Phi_i(\mathbf{x})
    = \sum_{i} \frac{q_i\,M_i}{4\pi}\,
      \phi(|\mathbf{x}-\mathbf{x}_i|).
\end{equation}
This approximation corresponds to the Newtonian limit and suffices
at large scales.

\subsubsection{Full nonlinear regime}

In regions where contributions from many sources overlap strongly
(dense matter clusters, stellar interiors, atomic nuclei), the field
equation becomes \textbf{nonlinear}:
\begin{equation}\label{eq:Phi-nonlinear}
    \mathcal{D}[\Phi]
    = \sum_{i} q_i \rho_i
    + \mathcal{N}[\Phi],
\end{equation}
where $\mathcal{N}[\Phi]$ describes the \textbf{self-interference} of
the field --- the nonlinear coupling of generated space with itself.
The explicit form of $\mathcal{N}[\Phi]$ is given in
Eq.~\eqref{eq:N-preview} above and in Definition~\ref{def:N}
(Section~\ref{sec:formalizm}); the full field equation is stated in
Theorem~\ref{thm:field-eq}.

The self-interference term $\mathcal{N}[\Phi]$ is responsible for:
\begin{itemize}[nosep]
  \item repulsion at intermediate scales (high spatial potential
        between nearby sources),
  \item a well at extremely small scales (confinement),
  \item amplification and cancellation of the field depending on the
        relative configuration of sources.
\end{itemize}

\begin{remark}[Interference, not mere superposition]
The word ``interference'' is deliberate.  In wave physics,
interference means that the result of superposition depends on
\emph{phase} --- constructive and destructive effects can occur.
In TGP the analogous role is played by the \textbf{field gradient}
$\nabla\Phi$: two nearby sources generate between them a region of
\emph{higher}~$\Phi$ than either produces alone (constructive
overlap); but the net force on objects depends on the
\emph{gradient} of~$\Phi$, which can point either toward or away
from the source, depending on the scale.  This is the mechanism
behind the three regimes (Section~\ref{sec:rezimy}).
\end{remark}

% ----------------------------------------------------------
\subsection{From $\Phi$ to the metric}\label{ssec:metryka}

The metric is not a fundamental object.  It is an \textbf{effective
description} of the configuration of~$\Phi$.  In regions where
$\Phi > 0$ and varies smoothly, one can define an effective metric:
\begin{equation}\label{eq:metric-from-Phi}
    g_{\mu\nu}^{\mathrm{eff}}
    = g_{\mu\nu}[\Phi, \partial\Phi, \ldots\,].
\end{equation}
The exact form of this relation is an \textbf{open problem} in TGP.
At the present stage one can state:
\begin{itemize}[nosep]
  \item the metric is defined only where $\Phi > 0$;
  \item a homogeneous $\Phi = \mathrm{const} > 0$ yields a flat
        (Minkowski) metric;
  \item gradients of $\Phi$ generate curvature --- this is gravity
        in TGP;
  \item large $\Phi$ with large gradients corresponds to a strong
        gravitational field.
\end{itemize}

\begin{hypothesis}[Minimal metric]\label{hyp:metric}
In the simplest ansatz the effective metric takes the form
\begin{equation}\label{eq:metric-minimal}
    ds^2 = -\frac{c(\Phi)^2}{c_0^2}\,dt^2
    + \frac{\Phi}{\PhiZero}\,\delta_{ij}\,dx^i dx^j,
\end{equation}
where $c(\Phi)$ is the local speed of light
(Section~\ref{sec:stale}) and $\PhiZero$ is the reference value.
The spatial part is modulated by $\Phi/\PhiZero$ --- more generated
space means a larger ``volume'' available in a given region.
\end{hypothesis}

This metric form is a \textbf{working hypothesis}, not an axiom.
Its consistency with observations must be verified.  A corrected
(exponential) form yielding correct post-Newtonian parameters is
derived in Section~\ref{ssec:metryka-popr}.
